\documentclass{article}

% Language setting
\usepackage[english]{babel}

% Set page size and margins
\usepackage[letterpaper,top=2cm,bottom=2cm,left=3cm,right=3cm,marginparwidth=1.75cm]{geometry}

% Useful packages
\usepackage{amsmath}
\usepackage{amssymb}
\usepackage{graphicx}
\usepackage[colorlinks=true, allcolors=black]{hyperref}
\usepackage{titlesec}
\usepackage{xparse}

\NewDocumentCommand{\qcq}{m o o m}{%
	\ensuremath{\forall #1 \IfValueT{#2}{, #2} \IfValueT{#3}{, #3}  \in  \mathbb{#4}}%
}
\NewDocumentCommand{\ext}{m o o m}{%
	\ensuremath{\exists #1 \IfValueT{#2}{, #2} \IfValueT{#3}{, #3}  \in  \mathbb{#4}}%
}

\title{\textbf{Équations Différentielles}}
\author{\textbf{Yassin Akkab}}

\newcommand{\R}{\mathbb{R}}
\newcommand{\N}{\mathbb{N}}
\newcommand{\Z}{\mathbb{Z}}
\newcommand{\C}{\mathbb{C}}
\newcommand{\Q}{\mathbb{Q}}
\newcommand{\D}{\mathbb{D}}

\begin{document}
	
	\maketitle
	\tableofcontents
	\newpage
	
	\section{Généralités}
	Une équation différentielle est une équation reliant une fonction inconnue (souvent notée $y$) et une ou plusieurs de ses dérivées (notées $y', y'', \dots$).
	\begin{itemize}
		\item Résoudre ou intégrer une équation différentielle consiste à déterminer toutes les fonctions définies et dérivables sur un intervalle $I$ qui vérifient l'équation. Ces fonctions sont appelées les \textbf{solutions} de l'équation.
		\item Si l'on impose des conditions aux limites ou des valeurs en un point (ex: $y(x_0) = y_0$), on parle de \textbf{problème de Cauchy} ou de conditions initiales.
	\end{itemize}
	
	\section{Équations différentielles du premier ordre $y' = ay + b$}
	Soient $a$ et $b$ deux nombres réels. On s'intéresse à l'équation :
	\[ (E_1) : y' = ay + b \]
	
	\subsection{Résolution générale}
	\textbf{Théorème} : Les solutions sur $\mathbb{R}$ de l'équation différentielle $(E_1)$ sont les fonctions de la forme :
	\begin{itemize}
		\item Si $a \ne 0$ :
		\[ y(x) = C e^{ax} - \frac{b}{a} \quad (\text{où } C \in \mathbb{R}) \]
		\item Si $a = 0$ :
		\[ y(x) = bx + C \quad (\text{où } C \in \mathbb{R}) \]
	\end{itemize}
	
	\subsection{Condition initiale}
	Pour tout couple de réels $(x_0, y_0)$, il existe une \textbf{unique} solution $f$ de l'équation $(E_1)$ vérifiant la condition initiale $f(x_0) = y_0$.
	
	\section{Équations différentielles du second ordre $y'' + ay' + by = 0$}
	Soient $a$ et $b$ deux nombres réels. On s'intéresse à l'équation linéaire homogène du second ordre à coefficients réels :
	\[ (E_2) : y'' + ay' + by = 0 \]
	
	\subsection{Équation caractéristique}
	On associe à l'équation $(E_2)$ l'équation algébrique du second degré appelée \textbf{équation caractéristique} :
	\[ (EC) : r^2 + ar + b = 0 \]
	Soit $\Delta = a^2 - 4b$ le discriminant de cette équation.
	
	\subsection{Résolution selon le signe de $\Delta$}
	\begin{enumerate}
		\item \textbf{Si $\Delta > 0$} :
		L'équation caractéristique $(EC)$ admet deux solutions réelles distinctes $r_1$ et $r_2$.
		Les solutions de l'équation $(E_2)$ sur $\mathbb{R}$ sont de la forme :
		\[ y(x) = \alpha e^{r_1 x} + \beta e^{r_2 x} \quad (\text{où } (\alpha, \beta) \in \mathbb{R}^2) \]
		
		\item \textbf{Si $\Delta = 0$} :
		L'équation caractéristique $(EC)$ admet une solution réelle double $r_0 = -\frac{a}{2}$.
		Les solutions de l'équation $(E_2)$ sur $\mathbb{R}$ sont de la forme :
		\[ y(x) = (\alpha x + \beta) e^{r_0 x} \quad (\text{où } (\alpha, \beta) \in \mathbb{R}^2) \]
		
		\item \textbf{Si $\Delta < 0$} :
		L'équation caractéristique $(EC)$ admet deux solutions complexes conjuguées $r_1 = p + iq$ et $r_2 = p - iq$ (où $p, q \in \mathbb{R}$ et $q \ne 0$).
		Les solutions de l'équation $(E_2)$ sur $\mathbb{R}$ sont de la forme :
		\[ y(x) = e^{px} \left( \alpha \cos(qx) + \beta \sin(qx) \right) \quad (\text{où } (\alpha, \beta) \in \mathbb{R}^2) \]
	\end{enumerate}
	
\end{document}
